\section{Приложение}

\subsection{Алфавитный индекс ключевых команд}
\begin{longtable}{>{\raggedright\arraybackslash}p{0.28\textwidth}>{\raggedright\arraybackslash}p{0.68\textwidth}}
\toprule
Команда & Назначение \\
\midrule
\texttt{arm()} & Включение моторов перед стартом \\
\texttt{disarm()} & Выключение моторов после посадки или в экстренной ситуации \\
\texttt{get\_battery\_status()} & Получение напряжения батареи \\
\texttt{get\_cv\_frame()} & Получение кадра камеры в формате OpenCV \\
\texttt{get\_dist\_sensor\_data()} & Чтение дальномера \\
\texttt{go\_to\_local\_point()} & Переход к точке в локальной системе координат \\
\texttt{land()} & Штатная посадка \\
\texttt{point\_reached()} & Проверка достижения целевой точки \\
\texttt{send\_rc\_channels()} & Отправка значений RC-каналов для ручного управления \\
\texttt{set\_manual\_speed()} & Управление линейными скоростями и рысканьем \\
\texttt{set\_manual\_speed\_body\_fixed()} & Управление скоростями в связанной системе координат \\
\texttt{takeoff()} & Автоматический взлёт \\
\bottomrule
\end{longtable}

\subsection{Карта скриптов для преподавателя}
\begin{table}[H]
\centering
\begin{tabularx}{\textwidth}{>{\raggedright\arraybackslash}p{0.39\textwidth}>{\raggedright\arraybackslash}X}
\toprule
Скрипт & Где использовать в курсе \\
\midrule
\texttt{mission\_examples/PiZero\_test.py} & Урок 2: первый взлёт и посадка \\
\texttt{mission\_examples/set\_manual\_speed.py} & Урок 4: прямолинейный полёт \\
\texttt{mission\_examples/go\_to\_local\_point.py} & Урок 5: полёт по точкам \\
\texttt{mission\_examples/circle\_flight.py} & Урок 6: базовые фигуры \\
\texttt{camera\_examples/camera\_calibration.py} & Урок 9: калибровка \\
\texttt{aruco\_examples/aruco\_flight.py} & Урок 11: визуальная обратная связь \\
\texttt{PioneerHumanTracking/human-tracking.py} & Урок 12: продвинутый CV \\
\texttt{group\_flight\_examples/swarm\_OPT.py} & Урок 13: роевые алгоритмы \\
\bottomrule
\end{tabularx}
\caption{Быстрый выбор скрипта под учебную цель}
\end{table}

\nocite{opencv_calibration}
\nocite{mediapipe_pose}
\nocite{python_threading}
